\documentclass[10pt]{article}
\usepackage[margin=0.8in]{geometry}
\usepackage{amsmath,microtype}
\newcommand{\guideheading}[1]{\par\medskip\noindent{\large\bfseries #1}\par\smallskip}
\begin{document}
\begin{center}
{\LARGE Guide: When Equal-Count Prefixes Cannot Merge}\par\medskip
Collatz proof project\quad---\quad September 5, 2026
\end{center}
\raggedright
\guideheading{What Lean proves}
Take any two natural starting numbers and run the shortcut Collatz map
for at most 31 steps. Suppose each path's multiplicative coefficient
stays at least one, the paths use the same number of odd steps, and they
end at the same value. Then the starting numbers were identical.
This holds for arbitrarily large seeds, not just a finite seed sample.

\guideheading{Why this matters for the current approach}
A smaller merging predecessor can help rule out a least noncontracting
seed. But the transport method needs a noncontracting replacement prefix.
The new result shows that keeping both the time and odd count unchanged
cannot provide such a distinct replacement within 31 steps. Different
lengths or odd counts remain available, as in the earlier merge rules.

\guideheading{The proof in two observations}
A noncontracting prefix of at least two steps must start with two odd
steps. Both seeds therefore have remainder three modulo four.

The exact endpoint formula consists of a multiplier times the starting
seed, plus a correction. A small checked table bounds that correction.
If two prefixes have the same time, odd count and endpoint, the bound
forces their starting numbers to differ by less than four. Two integers
with the same remainder modulo four cannot be distinct that close together.

\guideheading{Why the large search is not the proof}
An exploratory program found no equal-count collision through depth 30,
which contains 12,771,274 noncontracting prefixes. The Lean proof does
not trust that enumeration. Instead, the kernel checks a 32-by-32 table
of correction bounds and verifies the induction that applies it to every
natural seed. The resulting theorem reaches time 31.

\guideheading{What fails if we remove the hypotheses}
The seeds 31 and 95 both reach 182 after seven steps, and both prefixes
are noncontracting. They use six and five odd steps, respectively.
Lean checks this counterexample to dropping the equal-count condition.

The particular correction interval also fails at time 32. Lean certifies
a noncontracting prefix starting at 3,384,695,803 whose correction exceeds
that interval. This is not a counterexample to equal-count injectivity
at time 32. It means that extending the theorem needs a stronger argument.

\guideheading{What remains open}
The theorem is limited to 31 steps. No arbitrary-time injectivity result
has been proved, and neither nondivergence nor exclusion of all nontrivial
cycles follows. The full Collatz conjecture remains unproved.

\guideheading{Reproduce}
From \texttt{lean/}, run \texttt{lake build Collatz.NCPrefixInjective}.
The table generator and tests are in \texttt{analysis/}. The separate
\texttt{nc\_collision\_search.py} records a state-limit stop distinctly
from a completed search; it is not needed to check the Lean proof.
\end{document}
